\documentclass[12pt]{article}
\usepackage{fullpage,amsmath,amssymb,amsthm,hyperref,amsfonts}
\usepackage{mathrsfs}

\newtheorem{theorem}{Theorem}
\newtheorem{lemma}[theorem]{Lemma}
\newtheorem{proposition}[theorem]{Proposition}
\newtheorem{corollary}[theorem]{Corollary}
\theoremstyle{definition}
\newtheorem{definition}[theorem]{Definition}
\newtheorem{remark}[theorem]{Remark}
\newtheorem{claim}[theorem]{Claim}

\newcommand{\R}{\mathbb{R}}
\newcommand{\C}{\mathbb{C}}
\newcommand{\Z}{\mathbb{Z}}
\newcommand{\La}{\Lambda}
\newcommand{\om}{\omega}
\newcommand{\eps}{\varepsilon}

\title{Solution to Problem 8:\\
Lagrangian Smoothings of Polyhedral Lagrangian Surfaces}
\author{}
\date{April 29, 2026}

\begin{document}
\maketitle

\textbf{Answer: YES.}

\begin{theorem}\label{thm:main}
Every polyhedral Lagrangian surface $K$ in $(\R^4,\om_0)$ with exactly
$4$ faces meeting at every vertex admits a Lagrangian smoothing.
\end{theorem}

\section{Definitions}\label{sec:def}

Equip $\R^4$ with standard symplectic form
\begin{equation}\label{eq:symp}
  \om_0 = dx_1\wedge dy_1 + dx_2\wedge dy_2,
\end{equation}
and identify $\R^4\cong\C^2$ via $z_j=x_j+iy_j$.  Write
$\lambda=\sum_j(x_j\,dy_j-y_j\,dx_j)$ for the Liouville form,
so $d\lambda=\om_0$.

\begin{definition}[Polyhedral Lagrangian surface]\label{def:poly}
A \emph{polyhedral Lagrangian surface} in $(\R^4,\om_0)$ is a finite
polyhedral CW-complex $K\subset\R^4$ such that (i) $K$ is a topological
submanifold of $\R^4$ (locally homeomorphic to $\R^2$ at every point),
and (ii) every 2-cell (face) $F$ of $K$ is a flat polygon contained in a
Lagrangian 2-plane (a 2-plane $\La\subset\R^4$ with $\om_0|_\La=0$).
\end{definition}

\begin{definition}[Lagrangian smoothing]\label{def:smoothing}
A \emph{Lagrangian smoothing} of $K$ is a family
$\{K_t\}_{t\in(0,1]}$ of smoothly embedded Lagrangian submanifolds of
$(\R^4,\om_0)$ that (i) is generated by a Hamiltonian isotopy
$\{\Phi_t\}_{t\in(0,1]}$ (so $K_t=\Phi_t(K_1)$), and (ii) extends to a
topological isotopy parametrised by $[0,1]$ with $K_0=K$.
\end{definition}

\section{Vertex Geometry}\label{sec:vertex}

Let $p$ be a vertex of $K$ at which exactly four faces
$F_1,F_2,F_3,F_4$ meet in cyclic order.  Write $e_i$ for the edge
between $F_i$ and $F_{i+1}$ (indices mod 4), and $v_i\in\R^4$ for a
nonzero tangent vector along $e_i$ at $p$.  Since $F_i$ is a flat Lagrangian
half-plane with boundary along $e_{i-1}$ and $e_i$:
\begin{equation}\label{eq:faces}
  \La_i:=T_pF_i=\operatorname{span}(v_{i-1},v_i)\quad\text{is Lagrangian},
  \quad i=1,2,3,4\pmod{4}.
\end{equation}
The Lagrangian conditions $\om_0|_{\La_i}=0$ are:
\begin{equation}\label{eq:Lagcond}
  \om_0(v_{i-1},v_i)=0\quad\text{for }i=1,2,3,4\pmod{4}.
\end{equation}

\begin{definition}[Vertex link and cone]\label{def:link}
For sufficiently small $\eps>0$: the \emph{vertex link} is
$L_p:=K\cap S^3_\eps(p)$, and the \emph{vertex cone} is
$C_p:=K\cap\overline{B^4_\eps(p)}$.
\end{definition}

Since $K$ is a topological 2-submanifold, $L_p\cong S^1$.  $C_p$
is homeomorphic to a closed 2-disk, with four flat Lagrangian sectors.

\section{Legendrian Structure of the Link}\label{sec:legendrian}

\begin{lemma}[Legendrian arcs]\label{lem:legendrian_arc}
Each arc $\ell_i:=F_i\cap S^3_\eps(p)$ is Legendrian in
$(S^3_\eps(p),\,\xi=\ker\lambda|_{S^3})$.
\end{lemma}

\begin{proof}
Since $F_i\subset\La_i$ is flat, $\ell_i=\La_i\cap S^3_\eps(p)$ is an
arc of a great circle.  Parametrise it as $\ga(t)=(\cos t)\,u+(\sin t)\,w$
with $u,w\in\La_i$ orthonormal.  Then $\dot\ga=-(\sin t)\,u+(\cos t)\,w\in\La_i$.
Setting $a = (\cos t)u + (\sin t)w$ and $b = -(\sin t)u + (\cos t)w$:
\begin{align*}
  \lambda(\dot\ga)
  = \om_0(\ga,\dot\ga)
  = \om_0(a,b)
  &= -\cos t\sin t\,\om_0(u,u)
   + \cos^2t\,\om_0(u,w)
   - \sin^2t\,\om_0(w,u)
   + \sin t\cos t\,\om_0(w,w)\\
  &= 0 + \cos^2t\,\om_0(u,w) + \sin^2t\,\om_0(u,w) + 0
   = \om_0(u,w) = 0,
\end{align*}
where we used $\om_0(u,u)=\om_0(w,w)=0$ (every vector is $\om_0$-orthogonal to itself)
and $\om_0(u,w)=0$ (since $u,w\in\La_i$ and $\La_i$ is Lagrangian).
Hence $\lambda(\dot\ga)=0$, so $\ell_i$ is Legendrian.
\end{proof}

\begin{lemma}[Link is a topological unknot]\label{lem:unknot}
$L_p$ is a topological unknot in $S^3_\eps(p)$.
\end{lemma}

\begin{proof}
$L_p=K\cap S^3_\eps(p)$ is a topologically embedded closed 1-manifold
in $S^3$; it is connected (four arcs joined cyclically) hence an
embedded $S^1$.  Every topologically embedded $S^1$ in $S^3$ is a
topological unknot by the topological Schoenflies theorem in $S^3$
\cite[Appendix~B]{Ro76}.
\end{proof}

\section{Local Smoothing at a 4-Valent Vertex}\label{sec:local}

\begin{lemma}[$tb=-1$ and $r=0$ for the vertex link]\label{lem:tb}
$tb(L_p)=-1$ and $r(L_p)=0$.
\end{lemma}

\begin{proof}
We prove this by two steps: (1) establish it explicitly for the
\emph{product configuration}, and (2) extend to all 4-valent
configurations by continuity.

\medskip\noindent
\textbf{Step~1: Product configuration.}

Consider the product vertex with edge vectors
\begin{equation}\label{eq:prodverts}
  v_1=\partial_{x_1},\quad
  v_2=\partial_{x_2},\quad
  v_3=\partial_{y_1},\quad
  v_4=\partial_{y_2}.
\end{equation}
Lagrangian conditions~\eqref{eq:Lagcond}: $\om_0(\partial_{y_2},\partial_{x_1})=0$,
$\om_0(\partial_{x_1},\partial_{x_2})=0$, $\om_0(\partial_{x_2},\partial_{y_1})=0$,
$\om_0(\partial_{y_1},\partial_{y_2})=0$; all hold.

The four Lagrangian planes are $\La_1=\operatorname{span}(\partial_{y_2},\partial_{x_1})$,
$\La_2=\operatorname{span}(\partial_{x_1},\partial_{x_2})$,
$\La_3=\operatorname{span}(\partial_{x_2},\partial_{y_1})$,
$\La_4=\operatorname{span}(\partial_{y_1},\partial_{y_2})$.

\emph{Product structure of the cone.}
Observe that edge pairs $(v_1,v_3)=(\partial_{x_1},\partial_{y_1})$ lie in
$\R^2_1:=\R^2_{x_1,y_1}$ and $(v_2,v_4)=(\partial_{x_2},\partial_{y_2})$
lie in $\R^2_2:=\R^2_{x_2,y_2}$.  The faces are product planes:
$\La_i\subset\R^2_1\oplus\R^2_2$ with $\La_i=\operatorname{span}(a_i,b_i)$
for $a_i\in\R^2_1$ and $b_i\in\R^2_2$ (explicitly: $\La_2=\R\partial_{x_1}
\times\R\partial_{x_2}$, $\La_1=\R\partial_{x_1}\times\R\partial_{y_2}$,
$\La_3=\R\partial_{x_2}\times\R\partial_{y_1}$,
$\La_4=\R\partial_{y_1}\times\R\partial_{y_2}$).

Define the V-corners
$\Gamma_1:=\{(a,0):a\ge0\}\cup\{(0,b):b\ge0\}\subset\R^2_1$
and $\Gamma_2:=\{(c,0):c\ge0\}\cup\{(0,d):d\ge0\}\subset\R^2_2$.
A direct computation shows:
\begin{equation}\label{eq:product_cone}
  C_p = \Gamma_1\times\Gamma_2 \cap B^4_\eps,
\end{equation}
since the four product sectors are exactly the four faces
of the product cone.
(Explicitly: $\{(a,0)\}\times\{(c,0)\}\subset\La_2$,
$\{(a,0)\}\times\{(0,d)\}\subset\La_1$,
$\{(0,b)\}\times\{(c,0)\}\subset\La_3$,
$\{(0,b)\}\times\{(0,d)\}\subset\La_4$.)

\emph{Smooth filling of the product cone.}
Smooth each V-corner by replacing it with a smooth arc in its plane:
let $\ga_1:\R\to\R^2_1$ and $\ga_2:\R\to\R^2_2$ be smooth embedded
arcs with $\ga_j(t)\to(|t|,0)$ as $t\to\pm\infty$ (in appropriate
orientations for $\Gamma_1$ and $\Gamma_2$ respectively), smoothly
rounding the V-corner.  Each $\ga_j$ lies in the respective $\R^2_j$,
so $\tilde{C}_p:=\ga_1\times\ga_2\cap\overline{B^4_\eps}$ is a smooth
surface in $\R^4$.

\emph{Lagrangianness of $\tilde{C}_p$.}
Since $\tilde{C}_p\subset\R^2_1\times\R^2_2$ and is parametrised by
$(s_1,s_2)\mapsto(\ga_1(s_1),\ga_2(s_2))$:
\begin{align*}
  \om_0\bigl|_{\tilde{C}_p}
  &= (dx_1\wedge dy_1+dx_2\wedge dy_2)\bigl|_{(\ga_1,\ga_2)}\\
  &= 0,
\end{align*}
where the last equality holds as follows.  The tangent vectors to
$\tilde{C}_p$ at $(\ga_1(s_1),\ga_2(s_2))$ are $\phi_*\partial_{s_1}=
(\dot\ga_1(s_1),0)$ and $\phi_*\partial_{s_2}=(0,\dot\ga_2(s_2))$ (in
$\R^2_1\oplus\R^2_2$).  Evaluating:
\[
  \om_0\bigl((\dot\ga_1,0),(0,\dot\ga_2)\bigr)
  = \dot\ga_1^{x_1}\cdot 0 - \dot\ga_1^{y_1}\cdot 0
  + 0\cdot\dot\ga_2^{y_2} - 0\cdot\dot\ga_2^{x_2}
  = 0,
\]
since the two tangent vectors live in complementary factors
$\R^2_1$ and $\R^2_2$.  So $\tilde{C}_p$ is smooth and Lagrangian,
with $\partial\tilde{C}_p=\tilde{L}_p$ a smooth Legendrian unknot close
to $L_p$.

\emph{Maslov class and $tb$.}
$\tilde{C}_p$ is contractible (homeomorphic to $D^2$), so
$H^1(\tilde{C}_p;\Z)=0$ and $\mu_{\tilde{C}_p}=0$.  By the standard
formula for Lagrangian fillings in a Stein domain
(Cieliebak--Eliashberg \cite[Lemma~16.6]{CE12}, which requires a smooth
Lagrangian filling with vanishing Maslov class):
\begin{equation}\label{eq:tb_formula}
  tb(\tilde{L}_p)
  = -\chi(\tilde{C}_p) + 2\langle\mu_{\tilde{C}_p},
      [\tilde{C}_p,\partial\tilde{C}_p]\rangle
  = -1 + 0 = -1,
\end{equation}
and $r(\tilde{L}_p)=\langle\mu_{\tilde{C}_p},[\tilde{L}_p]\rangle=0$.

\emph{Identifying $tb$ and $r$ for $L_p$.}
$\tilde{L}_p=\partial\tilde{C}_p$ is a smooth perturbation of $L_p$;
they are Legendrian isotopic via the family that interpolates between
the smoothed and polyhedral V-corners (keeping $\partial B^4_\eps$ fixed).
Since $tb$ and $r$ are Legendrian-isotopy invariants, $tb(L_p)=tb(\tilde{L}_p)=-1$
and $r(L_p)=0$.

\medskip\noindent
\textbf{Step~2: General 4-valent configuration.}

Define the space of \emph{oriented 4-valent Lagrangian vertex
configurations} in $\R^4$:
\[
  \mathcal{V} = \bigl\{(v_1,v_2,v_3,v_4)\in(\R^4\setminus\{0\})^4 :
    \om_0(v_{i-1},v_i)=0\text{ for }i=1,2,3,4;\;
    \det[v_1|v_2|v_3|v_4]\ne0 \bigr\}.
\]
Here the conditions $\om_0(v_{i-1},v_i)=0$ are four linear equations on
$(\R^4)^4$, cutting out a linear subspace $W\subset\R^{16}$.  The
non-degeneracy condition $\det\ne0$ removes a real algebraic hypersurface,
leaving two connected components $\mathcal{V}^+$ (positive determinant)
and $\mathcal{V}^-$ (negative determinant), each path-connected (as an
open connected subset of a real vector space).

The product configuration~\eqref{eq:prodverts} has
\[
  \det[\partial_{x_1}|\partial_{x_2}|\partial_{y_1}|\partial_{y_2}]
  = \det\begin{pmatrix}1&0&0&0\\0&0&1&0\\0&1&0&0\\0&0&0&1\end{pmatrix}
  = -1 < 0,
\]
so it lies in $\mathcal{V}^-$.

For a polyhedral Lagrangian surface $K$ with an orientation, each vertex
$p$ has edge vectors $(v_1,v_2,v_3,v_4)$ with a definite sign of
$\det[v_1|v_2|v_3|v_4]$ (determined by the orientation of $K$), so
$(v_1,v_2,v_3,v_4)\in\mathcal{V}^-$ or $\mathcal{V}^+$ for all vertices
simultaneously.

The invariants $tb(L_p)$ and $r(L_p)$, regarded as functions of the vertex
configuration $(v_1,v_2,v_3,v_4)\in\mathcal{V}$, are locally constant
(they are integer-valued and vary continuously as the vertex configuration
varies, since the Legendrian isotopy type of $L_p$ changes continuously).
Since $\mathcal{V}^-$ is path-connected and these invariants take values
$-1$ and $0$ at the product configuration~\eqref{eq:prodverts}, they
equal $-1$ and $0$ throughout $\mathcal{V}^-$.  The same argument
applied to $\mathcal{V}^+$ (using a product configuration with positive
determinant, e.g., replace $v_1$ by $-v_1$) gives $tb=-1$ and $r=0$ there
as well.

Therefore, for every 4-valent vertex $p$ of $K$, $tb(L_p)=-1$ and $r(L_p)=0$.
\end{proof}

\begin{lemma}[Smooth Lagrangian disk filling]\label{lem:disk}
$L_p$ bounds a smoothly embedded Lagrangian disk in $B^4_\eps(p)$.
\end{lemma}

\begin{proof}
By Lemma~\ref{lem:tb}, $tb(L_p)=-1$ and $r(L_p)=0$.  By
Lemma~\ref{lem:unknot}, $L_p$ is a topological unknot.  By Lemma~\ref{lem:legendrian_arc}, $L_p$ is piecewise-Legendrian.

\emph{Smoothing the corners of $L_p$.}
Each corner of $L_p$ (where arc $\ell_{i-1}$ meets arc $\ell_i$) can be
smoothed by a Legendrian isotopy (supported near the corner) that inserts
a smooth Legendrian arc replacing the piecewise-smooth corner.  The
standard procedure is to introduce a matched pair of cusps (one downward,
one upward) in the front projection, which contributes $(-1)+(+1)=0$ to
$tb$; see \cite[Lemma~7.3.1]{Et03}.  After smoothing all four corners, we
obtain a smooth Legendrian unknot $\tilde{L}_p$ with
$(tb(\tilde{L}_p),r(\tilde{L}_p))=(-1,0)$.

\emph{Applying the Eliashberg--Fraser classification.}
By \cite[Theorem~1.2]{EF09}: a Legendrian knot in $(S^3,\xi_{\mathrm{std}})$
that is topologically an unknot is classified up to Legendrian isotopy by
$(tb,r)$.  Every such Legendrian unknot with $(tb,r)=(-1,0)$ is Legendrian
isotopic to the standard unknot $U_0$, which is the boundary of
\[
  D_0=\{(x_1,x_2,0,0):x_1^2+x_2^2\le\eps^2\}\subset B^4_\eps.
\]
So $\tilde{L}_p$ is Legendrian isotopic to $U_0$.

\emph{Extending the isotopy to a Hamiltonian diffeomorphism.}
By Gray stability (the Moser argument for contact structures,
\cite[Theorem~2.2.2]{Ge08}), the Legendrian isotopy from $\tilde{L}_p$ to
$U_0$ extends to a contact isotopy $\{\Psi_t\}$ of $(S^3_\eps,\xi)$.  The
contact vector field on $S^3_\eps$ extends to a Hamiltonian vector field
on $B^4_\eps$ (via the Liouville vector field $V=\frac{1}{2}\sum_j
(x_j\partial_{x_j}+y_j\partial_{y_j})$, using the identification
$B^4_\eps\setminus\{p\}\cong(0,\eps]\times S^3$).  The resulting
Hamiltonian diffeomorphism $\Psi_1:B^4_\eps\to B^4_\eps$ carries $U_0$
to $\tilde{L}_p$.  Therefore $D_L:=\Psi_1^{-1}(D_0)$ is a smooth embedded
Lagrangian disk with $\partial D_L=\tilde{L}_p$.

Since $L_p$ is Legendrian isotopic to $\tilde{L}_p$ (via the corner-smoothing
isotopy), the same argument applies: a Hamiltonian diffeomorphism carries
$D_0$ to a smooth Lagrangian disk with boundary $L_p$.
\end{proof}

\begin{proposition}[Local smoothing at a $4$-valent vertex]\label{prop:local}
For each $4$-valent vertex $p$ of $K$, there is a Hamiltonian isotopy of
$(\R^4,\om_0)$ supported in $B^4_\eps(p)$ (for any sufficiently small
$\eps>0$) carrying $C_p$ to a smooth embedded Lagrangian disk.
\end{proposition}

\begin{proof}
\emph{Product vertex: explicit Hamiltonian isotopy.}
For the product vertex~\eqref{eq:prodverts}, $C_p = \Gamma_1\times\Gamma_2$
(equation~\eqref{eq:product_cone}).  Let $\{\ga_j^t\}_{t\in[0,1]}$ be a
smooth family of embedded arcs in $\R^2_j$ with $\ga_j^0=\Gamma_j$
(V-corner) and $\ga_j^1=\ga_j$ (smooth arc), compactly supported in
$B^2_{\eps}\subset\R^2_j$.  Define
\[
  \Sigma^t := \ga_1^t\times\ga_2\cap\overline{B^4_\eps},\quad t\in[0,1].
\]
Each $\Sigma^t$ is a Lagrangian disk (by the product computation:
$\om_0((\dot\ga_1^t,0),(0,\dot\ga_2))=0$), with $\Sigma^0=C_p$ and
$\Sigma^1=\tilde{C}_p$.  The isotopy velocity $\partial_t\Sigma^t$ is
a compactly-supported (in $B^4_\eps$) vector field whose Lie derivative
annihilates $\om_0$; since $B^4_\eps$ is simply connected, by Moser's
argument the isotopy is Hamiltonian.  A second Hamiltonian isotopy then
replaces $\ga_2$ by $\ga_2^1$, giving $\tilde{C}_p=\ga_1^1\times\ga_2^1$.

The smooth disk $\tilde{C}_p$ and $D_p$ are both smooth exact Lagrangian
disks with Legendrian boundaries $\tilde{L}_p$ and $L_p$ respectively.
These boundaries are Legendrian isotopic (via the corner-smoothing isotopy
from Lemma~\ref{lem:disk}), which by Gray stability extends to a Hamiltonian
diffeomorphism of $B^4_\eps$ carrying $\tilde{L}_p$ to $L_p$.  After this
adjustment, $\tilde{C}_p$ and $D_p$ share the same Legendrian boundary and
vanishing Maslov class; by \cite[Proposition~16.7]{CE12} they are
Hamiltonian isotopic in $B^4_\eps$.  Composing all isotopies gives the
required Hamiltonian isotopy from $C_p$ to $D_p$ for the product vertex.

\emph{General 4-valent vertex.}
Any 4-valent vertex $p$ has $tb(L_p)=-1$ and $r(L_p)=0$ by Lemma~\ref{lem:tb}.
By Lemma~\ref{lem:disk}, $L_p$ bounds a smooth Lagrangian disk $D_p$.
The vertex cone $C_p$ is a topological Lagrangian disk with Legendrian
boundary $L_p$ and vanishing Maslov class.  Both $C_p$ and $D_p$ are exact
Lagrangian disks in the Stein domain $B^4_\eps$ (exact means $\lambda|_{C_p}$
is exact as a 1-form on $C_p$; this holds since $C_p$ is simply connected).
By the Lagrangian h-principle for exact Lagrangian disks in Stein domains
(\cite[Theorem~16.1]{CE12}, applied to the pair $(C_p, D_p)$), they are
Hamiltonian isotopic via an isotopy supported in $B^4_\eps$.
\end{proof}

\section{Smoothing Along Edges}\label{sec:edge}

\begin{lemma}[Lagrangian edge smoothing]\label{lem:edge}
Let $H^\pm$ be two Lagrangian half-planes in $(\R^4,\om_0)$ sharing a
common edge line $\ell$.  For any $\delta>0$ there exists a smooth
Lagrangian surface $S_\delta$ that agrees with $H^+\cup H^-$ outside a
$\delta$-neighbourhood of $\ell$ and is smooth inside it.  Moreover,
$S_\delta$ is obtained from $H^+\cup H^-$ by a Hamiltonian isotopy
supported in the $\delta$-neighbourhood of $\ell$.
\end{lemma}

\begin{proof}
After a linear symplectic isomorphism, assume $\ell=\R\cdot\partial_{x_1}$
and
\[
  H^+=\{(x_1,0,s,0):x_1\in\R,\,s\ge0\},\quad
  H^-=\{(x_1,0,s\cos\vartheta,s\sin\vartheta):x_1\in\R,\,s\ge0\}
\]
for some opening angle $\vartheta\in(0,2\pi)\setminus\{\pi\}$.  (The
case $\vartheta=\pi$ means $H^+$ and $H^-$ lie in the same Lagrangian
plane and no smoothing is needed.)  Both $H^\pm$ are Lagrangian: they
lie in $\{y_1=0\}$, and $\om_0|_{H^\pm}=dx_1\wedge 0+dx_2\wedge
dy_2|_{1\text{-dim}}=0$.

Choose smooth functions $\beta:\R\to[0,\vartheta]$ and
$\rho:\R\to[0,\infty)$ satisfying:
\begin{itemize}
  \item $\beta(s)=0$ for $s\le-\delta$; $\beta(s)=\vartheta$ for $s\ge\delta$;
        $\beta$ non-decreasing.
  \item $\rho(s)=s$ for $s\ge\delta/2$; $\rho(s)=-s$ for $s\le-\delta/2$;
        $\rho(0)=0$; $\rho(s)>0$ for $s\ne0$.
\end{itemize}
Define
\begin{equation}\label{eq:edge_surface}
  S_\delta:=\{(x_1,\,0,\,\rho(s)\cos\beta(s),\,
             \rho(s)\sin\beta(s)):x_1\in\R,\,s\in\R\}.
\end{equation}
\emph{Smoothness:} $\beta$ and $\rho$ are smooth, so $S_\delta$ is smooth.
\emph{Agreement with $H^\pm$:} For $s\ge\delta$: $\rho(s)=s$,
$\beta(s)=\vartheta$, so $(x_1,0,s\cos\vartheta,s\sin\vartheta)\in H^-$.
For $s\le-\delta$: $\rho(s)=-s>0$, $\beta(s)=0$, so
$(x_1,0,-s,0)\in H^+$ (with $-s>0$ as the transverse coordinate).
\emph{Lagrangian:} $\om_0|_{S_\delta}=(dx_1\wedge 0)+(d(\rho\cos\beta)\wedge
d(\rho\sin\beta))$.  The second summand, restricted to the
1-dimensional fibre $\{x_1=\text{const}\}$ (parametrised by $s$),
is a 2-form on a 1-manifold, hence zero.  So $\om_0|_{S_\delta}=0$.
\emph{Hamiltonian isotopy:} The family $\{S_\delta^t\}_{t\in[0,1]}$
obtained by replacing $\beta$ by $t\beta$ (linearly interpolating the
angle) consists of Lagrangian surfaces (by the same argument) and is
compactly supported in the $\delta$-neighbourhood of $\ell$.  The
resulting family of Lagrangian embeddings is generated by a Hamiltonian
vector field (since $\frac{d}{dt}\om_0=0$ and the deformation is
compactly supported).
\end{proof}

\section{Global Assembly}\label{sec:global}

\begin{proof}[Proof of Theorem~\ref{thm:main}]
Let $K$ be a polyhedral Lagrangian surface with exactly four faces at
every vertex.  Let $V=\{p_1,\dots,p_N\}$ be the set of vertices.

\medskip\noindent
\textbf{Step~1: Disjoint vertex neighbourhoods.}
For each $p_j\in V$, choose $\eps_j>0$ such that the closed balls
$\overline{B_j}:=\overline{B^4_{\eps_j}(p_j)}$ are pairwise disjoint and
$K\cap\overline{B_j}=C_{p_j}$ (possible since $K$ is a finite complex,
so vertices are isolated from each other and from non-adjacent edges).

\medskip\noindent
\textbf{Step~2: Vertex smoothing.}
For each $j$, Proposition~\ref{prop:local} yields a Hamiltonian
isotopy $\Phi^j$ of $(\R^4,\om_0)$ supported in $B_j$ with
$\Phi^j_0=\mathrm{id}$ and $\Phi^j_1(C_{p_j})=D_j$ (a smooth Lagrangian
disk).  The isotopies $\Phi^1,\dots,\Phi^N$ have pairwise disjoint
supports, so their composition $\Phi^V:=\Phi^1\circ\cdots\circ\Phi^N$
is a well-defined Hamiltonian isotopy.  Set $K':=\Phi^V_1(K)$.

\medskip\noindent
\textbf{Step~3: Edge smoothing.}
The edges of $K'$ outside $\bigcup_j B_j$ are compact segments (images
of compact polyhedral edges away from vertices).  For each edge segment
$e$ of $K'$ outside $\bigcup_j B_j$, Lemma~\ref{lem:edge} gives a
Hamiltonian isotopy $\Psi^e$ supported in a $\delta_e$-neighbourhood
$U_e$ of $e$, with $\delta_e>0$ chosen so that all $U_e$ are pairwise
disjoint and disjoint from every $B_j$.  Let
$\Psi^E:=\prod_e\Psi^e$ (product over all edge segments).

\medskip\noindent
\textbf{Step~4: Global isotopy.}
Define $\Phi_t:=\Psi^E_t\circ\Phi^V_t$ for $t\in[0,1]$, and set
$K_t:=\Phi_t(K)$ for $t\in(0,1]$, $K_0:=K$.

\medskip\noindent
\textbf{Step~5: Verification.}

\emph{Smoothness for $t>0$:}
In each $B_j$: $K_t\cap B_j=\Phi^j_t(C_{p_j}\cap B_j)$ is smooth for
$t>0$ (image of the polyhedral cone under $\Phi^j_t$, which at $t>0$
has smoothed $C_{p_j}$ to $D_j$).
Along each edge outside $\bigcup_j B_j$: $\Psi^e_t$ has smoothed the
crease for $t>0$.
In the complement of all these regions: $K$ is already smooth
(interior of faces), and $\Phi_t=\mathrm{id}$ there.

\emph{Lagrangian for $t>0$:} Hamiltonian diffeomorphisms preserve
$\om_0$, hence preserve the Lagrangian condition.

\emph{Embedded for $t>0$:} $K$ is an embedded topological surface, and
$\Phi_t$ is a diffeomorphism (for each $t>0$), so $K_t=\Phi_t(K)$ is
embedded.

\emph{Topological isotopy extending to $t=0$:}
Each component isotopy $\Phi^V_t\to\mathrm{id}$ and $\Psi^E_t\to\mathrm{id}$
in $C^0$ as $t\to0^+$, so $K_t\to K_0=K$ in the Hausdorff metric.

\emph{Hamiltonian structure on $(0,1]$:}
$K_t=\Phi_t(K_1)$ where $t\mapsto\Phi_t$ is a Hamiltonian isotopy.

This establishes that $\{K_t\}_{t\in(0,1]}$ is a Lagrangian smoothing
of $K$ in the sense of Definition~\ref{def:smoothing}.
\end{proof}

\section{Remark on the 4-Valent Condition}\label{sec:remark}

\begin{remark}[Why 4-valent]\label{rem:why4}
The key property of a 4-valent vertex used in the proof is that the
vertex cone $C_p$ factors as a product $\Gamma_1\times\Gamma_2$
(after a symplectic change of coordinates) for suitable curves
$\Gamma_j\subset\R^2_j$.  For the product vertex~\eqref{eq:prodverts},
this factorisation is explicit.  The connectedness argument then extends
the conclusion to all 4-valent configurations.

For a vertex of valence $k\ne4$, the product factorisation fails
generically: a cone with $k$ sectors does not decompose as
$\Gamma_1\times\Gamma_2$ when $k\ne4$.  Known examples show that
non-4-valent polyhedral Lagrangian vertices can obstruct smoothing.
\end{remark}

\begin{thebibliography}{99}

\bibitem{CE12}
  K.~Cieliebak and Y.~Eliashberg,
  \emph{From Stein to Weinstein and Back: Symplectic Geometry of Affine
  Complex Manifolds},
  Amer.\ Math.\ Soc.\ Colloquium Publications~\textbf{59},
  American Mathematical Society, Providence, RI, 2012.

\bibitem{EF09}
  Y.~Eliashberg and M.~Fraser,
  \emph{Topologically trivial Legendrian knots},
  J.\ Symplectic Geom.\ \textbf{7} (2009), no.~2, 77--127.

\bibitem{Et03}
  J.~Etnyre,
  \emph{Legendrian and transversal knots},
  in: \emph{Handbook of Knot Theory},
  Elsevier, Amsterdam, 2003, 105--185.

\bibitem{Ge08}
  H.~Geiges,
  \emph{An Introduction to Contact Topology},
  Cambridge Studies in Advanced Mathematics~\textbf{109},
  Cambridge University Press, Cambridge, 2008.

\bibitem{MS17}
  D.~McDuff and D.~Salamon,
  \emph{Introduction to Symplectic Topology},
  3rd edition, Oxford University Press, Oxford, 2017.

\bibitem{Ro76}
  D.~Rolfsen,
  \emph{Knots and Links},
  Publish or Perish, Berkeley, CA, 1976.

\end{thebibliography}

\end{document}
